\chapter{Results, Comparative Evaluation, and Discussion}
\label{ch:results}

\section{Introduction}
\label{sec:res_intro}
This chapter presents the empirical results and scientific findings of TerraShift AI across both standard academic benchmarks (OSCD) and our authenticated Saudi Deployment Multi-Model Benchmark (v6.1). All evaluations strictly prohibit synthetic benchmark fallbacks (\texttt{allow\_synthetic=False}), ingesting genuine Copernicus Sentinel-2 Level-2A surface reflectance via AWS Earth Search STAC and authentic Google AlphaEarth V1 64-D embeddings (\texttt{GOOGLE/SATELLITE\_EMBEDDING/V1/ANNUAL}) alongside Google Dynamic World V1 (\texttt{GOOGLE/DYNAMICWORLD/V1}) from Google Earth Engine (\texttt{project='terrashift'}). We report both **Pooled (micro)** and **Corridor Macro-Averaged** metrics side-by-side under our Localization-Aware 1-to-1 Bipartite Event-Matching Protocol.

\section{OSCD Academic Benchmark Reproduction}
\label{sec:res_oscd}
To establish an empirical baseline against computer vision literature, we evaluated four baseline model configurations against the 10 withheld evaluation scenes of the Onera Satellite Change Detection (OSCD) dataset \citep{daudt2018urban}. Results are reported across 3,077,936 evaluated ground-truth pixels in Table \ref{tab:oscd_benchmark_results}.

\begin{table}[htbp]
\centering
\caption{Quantitative performance comparison on the OSCD 10-scene test set.}
\label{tab:oscd_benchmark_results}
\begin{tabularx}{\textwidth}{@{}lXXXXl@{}}
\toprule
\textbf{Model Configuration} & \textbf{Precision} & \textbf{Recall} & \textbf{F1 Score} & \textbf{IoU Score} & \textbf{Pixel FPR} \\
\midrule
M0: Untrained Baseline (Random Weights) & 0.0542 & 0.4812 & 0.0980 & 0.0520 & 18.42\% \\
Statistical Spectral Baseline (Median/MAD) & 0.3210 & 0.4480 & 0.3740 & 0.2301 & 4.12\% \\
M1: Supervised Siamese U-Net (BCE Loss) & 0.3640 & 0.3980 & 0.3800 & 0.2350 & 3.05\% \\
\textbf{M1: Siamese U-Net + Calibrated Fusion} & \textbf{0.3780} & \textbf{0.4520} & \textbf{0.4117} & \textbf{0.2591} & \textbf{2.01\%} \\
\midrule
\textit{Reference Literature: FC-Siam-diff \citep{daudt2018fully}} & \textit{0.4680} & \textit{0.5110} & \textit{0.4886} & \textit{0.3232} & \textit{1.89\%} \\
\bottomrule
\end{tabularx}
\end{table}

\subsection{Analysis of Academic Baseline Reproduction}
The calibrated Siamese U-Net achieved an F1 score of **0.4117** and an IoU of **0.2591**, representing an 8.3\% relative improvement over the uncalibrated supervised baseline (0.3800) and outperforming the statistical spectral baseline (0.3740). The pixel-level False Positive Rate was constrained to 2.01\%.

However, this academic experiment revealed a fundamental limitation: training exclusively on RGB channels forces the neural network to approximate complex surface scattering from three narrow visible bands, lacking the diagnostic power of short-wave infrared (SWIR) and red-edge reflectance. When deployed to desert environments, RGB models suffer catastrophic false positive rates, demonstrating that higher OSCD scores on European imagery do not translate into operational reliability in the Middle East.

\section{The Saudi Multi-Model Benchmark Leaderboard (v6.1)}
\label{sec:res_saudi_leaderboard}
We conducted a comprehensive head-to-head empirical evaluation across all four Saudi reference corridors ($91.19\text{ km}^2$), evaluating all seven candidate model paradigms under our Localization-Aware 1-to-1 Event-Matching Protocol ($\text{PoC} \ge 0.15, \text{IoU} \ge 0.05$). Table \ref{tab:saudi_multi_model_leaderboard} presents both **Pooled (micro)** and **Corridor Macro-Averaged** F1, Precision, and Recall side-by-side.

\begin{table}[htbp]
\centering
\caption{Saudi Deployment Multi-Model Benchmark Leaderboard (v6.1 Authenticated GEE + STAC across 4 Corridors, $91.19\text{ km}^2$, 19 Reference Events).}
\label{tab:saudi_multi_model_leaderboard}
\resizebox{\textwidth}{!}{%
\begin{tabular}{@{}cllllllll@{}}
\toprule
\textbf{Rank} & \textbf{Model Paradigm} & \textbf{Pooled F1} & \textbf{Macro F1} & \textbf{Pooled / Macro Prec} & \textbf{Recall} & \textbf{False Alarms} & \textbf{Neg Hits} & \textbf{CPU Latency} \\
\midrule
\textbf{\#1} & \textbf{TerraShift v6.0 Hybrid (Siamese + AlphaEarth + L2A)} & \textbf{79.2\%} & \textbf{82.5\%} & \textbf{65.5\% / 70.1\%} & \textbf{100.0\%} & \textbf{10 (11.0/100km$^2$)} & \textbf{0 / 4} & \textbf{1.485 s} \\
\textbf{\#2} & \textbf{TerraShift v5.1 (Tuned Siamese + L2A Spectral Guards)} & \textbf{53.5\%} & \textbf{64.1\%} & \textbf{36.5\% / 47.2\%} & \textbf{100.0\%} & \textbf{33 (36.2/100km$^2$)} & \textbf{0 / 4} & \textbf{1.450 s} \\
\#3 & Google AlphaEarth V1 64-D Cosine (Model A, $\tau=0.12$) & 27.0\% & 31.4\% & 15.9\% / 19.1\% & 89.5\% & 90 (98.7/100km$^2$) & 23 / 4 & 0.498 s \\
\#4 & Pure Multispectral Physical Baseline ($\Delta\text{NDBI}/\Delta\text{NDVI}$) & 26.5\% & 26.7\% & 15.4\% / 15.5\% & 94.7\% & 99 (108.6/100km$^2$) & 0 / 4 & 0.556 s \\
\#5 & Untuned Siamese U-Net Baseline (RGB) & 23.8\% & 25.3\% & 13.6\% / 14.6\% & 94.7\% & 114 (125.0/100km$^2$) & 5 / 4 & 0.518 s \\
\#6 & Google Dynamic World V1 LULC Transition & 19.7\% & 27.9\% & 14.3\% / 33.5\% & 31.6\% & 36 (39.5/100km$^2$) & 6 / 4 & 0.497 s \\
\#7 & AlphaEarth + Dynamic World Transition (Model B) & 16.3\% & 25.9\% & 13.3\% / 33.6\% & 21.1\% & 26 (28.5/100km$^2$) & 4 / 4 & 0.464 s \\
\bottomrule
\end{tabular}%
}
\end{table}

\subsection{Key Scientific Findings from the Verified Leaderboard}
The empirical results in Table \ref{tab:saudi_multi_model_leaderboard} provide decisive, audit-verified answers to our central research question:

\begin{enumerate}
    \item \textbf{Decisive Superiority of TerraShift v6.0 Hybrid (Rank \#1):} By combining bitemporal Siamese U-Net spatial parcel delineation with Google AlphaEarth V1 64-D embedding consensus and Sentinel-2 L2A physical spectral guards ($\Delta\text{NDBI}$, $\Delta\text{NDVI}$), TerraShift v6.0 Hybrid achieves **79.2\% Pooled Event F1** and **82.5\% Macro Event F1**, with **65.5\% Pooled Strict Precision** (**70.1\% Macro Precision**; **76.0\% Pooled Validity Precision**) and **100.0\% Recall** (detecting all 19/19 reference parcels across all 4 corridors). It cuts total false alarms by **91.2\%** (from 114 down to 10) and achieves **0 Negative Control Zone hits**.
    \item \textbf{High-Cadence Tactical Autonomy of TerraShift v5.1 (Rank \#2):} Operating purely on 5-day Sentinel-2 L2A STAC imagery with zero GEE dependency, TerraShift v5.1 achieves **53.5\% Pooled Event F1** (**64.1\% Macro F1**), **36.5\% Pooled Precision** (**47.2\% Macro Precision**), and **100.0\% Recall** (19/19), reducing false alarms by **71.1\%** (114 down to 33) with **0 Negative Control Zone hits**.
    \item \textbf{Why Standalone Google AlphaEarth V1 Fails in Hyper-Arid Desert (Rank \#3):} When evaluated on authentic Google Earth Engine embeddings without physical spectral guards, standalone AlphaEarth V1 cosine dissimilarity ($\tau=0.12$) achieves only **27.0\% Pooled F1** (**31.4\% Macro F1**) and generates **90 false alarms** across $91.19\text{ km}^2$—including **23 false alarms inside undisturbed Negative Control Zones** (3 in Riyadh limestone plateau, 9 in Jeddah coastal sabkha, 1 in Al-Jouf dunes, and 10 in Al-Kharj desert regolith). Because AlphaEarth encodes general annual spatio-temporal appearance, inter-annual sand migration, moisture evaporation on sabkhas, and BRDF shifts induce cosine distances $d_{\cos} \in [0.12, 0.22]$ over invariant desert terrain.
    \item \textbf{Why Standalone Google Dynamic World V1 Fails on Arid Earthworks (Ranks \#6 \& \#7):} Standalone Google Dynamic World V1 LULC transition achieves only **19.7\% Pooled F1** (**31.6\% Recall**, detecting only 6 of 19 events) and triggers **6 Negative Control Zone hits**. Because bright limestone grading pads and unpaved desert roads share almost identical visible reflectance with bare desert regolith, Dynamic World's global classifier labels both epochs as \textit{Bare Ground}, missing 68.4\% of real Saudi construction and infrastructure events.
\end{enumerate}

\section{In-Depth Corridor Case Studies}
\label{sec:res_case_studies}
Table \ref{tab:per_corridor_results} breaks down the per-corridor performance of **TerraShift v6.0 Hybrid** and **TerraShift v5.1** side-by-side across all four Saudi evaluation corridors.

\begin{table}[htbp]
\centering
\caption{Per-Corridor Performance Breakdown: TerraShift v6.0 Hybrid vs. TerraShift v5.1 ($91.19\text{ km}^2$).}
\label{tab:per_corridor_results}
\resizebox{\textwidth}{!}{%
\begin{tabular}{@{}llllllllll@{}}
\toprule
\textbf{Corridor Name} & \textbf{Area ($\text{km}^2$)} & \textbf{v6.0 Alerts (TP/FP)} & \textbf{v6.0 Prec} & \textbf{v6.0 F1} & \textbf{v5.1 Alerts (TP/FP)} & \textbf{v5.1 Prec} & \textbf{v5.1 F1} & \textbf{Recall} & \textbf{Neg Hits} \\
\midrule
Riyadh ROSHN Sedra & 17.88 & 7 (5 TP / 2 FP) & \textbf{71.4\%} & \textbf{83.3\%} & 10 (5 TP / 5 FP) & 50.0\% & 66.7\% & \textbf{100.0\%} (5/5) & 0 / 1 \\
Jeddah North Obhur & 18.31 & 4 (4 TP / 0 FP) & \textbf{100.0\%} & \textbf{100.0\%} & 5 (4 TP / 1 FP) & 80.0\% & 88.9\% & \textbf{100.0\%} (4/4) & 0 / 1 \\
Al-Jouf Wadi As-Sirhan & 26.82 & 8 (5 TP / 3 FP) & \textbf{62.5\%} & \textbf{76.9\%} & 18 (5 TP / 13 FP) & 27.8\% & 43.5\% & \textbf{100.0\%} (5/5) & 0 / 1 \\
\textbf{Al-Kharj (Held-Out)} & \textbf{28.18} & \textbf{8 (5 TP / 3 FP)} & \textbf{62.5\%} & \textbf{76.9\%} & \textbf{16 (5 TP / 11 FP)} & \textbf{31.2\%} & \textbf{47.6\%} & \textbf{100.0\%} (5/5) & \textbf{0 / 1} \\
\midrule
\textbf{Pooled Total (Micro)} & \textbf{91.19} & \textbf{29 (19 TP / 10 FP)} & \textbf{65.5\%} & \textbf{79.2\%} & \textbf{52 (19 TP / 33 FP)} & \textbf{36.5\%} & \textbf{53.5\%} & \textbf{100.0\%} (19/19) & \textbf{0 / 4} \\
\textbf{Corridor Macro Mean} & \textbf{91.19} & \textbf{—} & \textbf{70.1\%} & \textbf{82.5\%} & \textbf{—} & \textbf{47.2\%} & \textbf{64.1\%} & \textbf{100.0\%} (19/19) & \textbf{0 / 4} \\
\bottomrule
\end{tabular}%
}
\end{table}

\subsection{Case Study 1: Riyadh ROSHN Sedra (Limestone Plateau Grading)}
In Riyadh Sedra ($17.88\text{ km}^2$), the untuned RGB baseline generated 39 alerts (34 false alarms, 12.8\% precision). TerraShift v5.1 reduced alerts to 10 (5 TP, 5 FP, **50.0\% precision**, **66.7\% F1**), while **TerraShift v6.0 Hybrid** further suppressed desert limestone false alarms to achieve **71.4\% strict precision** and **83.3\% F1** (5 TP, only 2 FP; **83.3\% validity precision**), with 0 hits in the limestone plateau negative control zone (where standalone AlphaEarth triggered 3 false alarms).

\subsection{Case Study 2: Jeddah North Obhur (Coastal Infill and Sabkha)}
Jeddah Obhur ($18.31\text{ km}^2$) achieved the highest precision across the national benchmark. Whereas standalone AlphaEarth generated 18 false alarms (including 9 false alarms inside the coastal sabkha negative control zone due to inter-annual brine moisture shifts), **TerraShift v6.0 Hybrid** achieved **100.0\% strict precision** and **100.0\% F1** (4 alerts, 4 TPs, **0 false alarms**), and **TerraShift v5.1** achieved **80.0\% strict precision** (**88.9\% F1**, 4 TPs, 1 FP).

\subsection{Case Study 3: Al-Jouf Wadi As-Sirhan (Center-Pivot Agriculture)}
In Al-Jouf ($26.82\text{ km}^2$), the system verified 50-hectare circular center-pivot irrigation fields. Standalone AlphaEarth generated 37 false alarms across fallow desert periphery circles (11.9\% precision). TerraShift v5.1 achieved **100.0\% recall** with 13 false alarms (27.8\% precision), while **TerraShift v6.0 Hybrid** combined AlphaEarth representation shifts with directional $\Delta\text{NDVI}$ guards to achieve **62.5\% strict precision** and **76.9\% F1** (5 TPs, 3 FPs, **83.3\% validity precision**) with 0 dune negative control hits.

\subsection{Case Study 4: Held-Out Al-Kharj Generalization Test}
Evaluated on the completely held-out Al-Kharj corridor ($28.18\text{ km}^2$) under frozen hyperparameters, **TerraShift v6.0 Hybrid** detected all 5 reference events (**100.0\% recall**, **62.5\% strict precision**, **76.9\% F1**, **71.4\% validity precision**), and **TerraShift v5.1** detected all 5 reference events (**100.0\% recall**, **31.2\% strict precision**, **47.6\% F1**), with 0 false alarms in the southern desert negative control zone (compared to 10 negative control zone hits by standalone AlphaEarth).

\section{Ablation & Localization Sensitivity Analysis}
\label{sec:res_ablation}
Table \ref{tab:ablation_comparison} isolates the progressive improvement from the Untuned RGB Siamese baseline to TerraShift v5.1 (Physical Spectral Guards) and TerraShift v6.0 Hybrid (Siamese + AlphaEarth + Physical Guards), reporting both Pooled and Macro metrics.

\begin{table}[htbp]
\centering
\caption{Progressive Ablation Analysis across All 4 Corridors ($91.19\text{ km}^2$, 19 Reference Parcels).}
\label{tab:ablation_comparison}
\begin{tabularx}{\textwidth}{@{}lXXX@{}}
\toprule
\textbf{Evaluation Metric} & \textbf{Untuned RGB (M0)} & \textbf{TerraShift v5.1 (L2A)} & \textbf{TerraShift v6.0 Hybrid} \\
\midrule
Total Predicted Alerts & 132 alerts & 52 alerts (-60.6\%) & \textbf{29 alerts (-78.0\%)} \\
Total False Positive Alerts & 114 alerts & 33 alerts (-71.1\%) & \textbf{10 alerts (-91.2\%)} \\
Pooled Strict 1-to-1 Precision & 13.6\% & 36.5\% (+22.9 pp) & \textbf{65.5\% (+51.9 pp)} \\
Corridor Macro Strict Precision & 14.6\% & 47.2\% (+32.6 pp) & \textbf{70.1\% (+55.5 pp)} \\
Pooled Event Validity Precision & 15.8\% & 45.2\% (+29.4 pp) & \textbf{76.0\% (+60.2 pp)} \\
Event Recall (19 Reference Parcels) & 94.7\% (18/19) & 100.0\% (19/19) & \textbf{100.0\% (19/19)} \\
Pooled Event F1 Score (Micro) & 23.8\% & 53.5\% (+29.7 pp) & \textbf{79.2\% (+55.4 pp)} \\
Corridor Macro Event F1 Score & 25.3\% & 64.1\% (+38.8 pp) & \textbf{82.5\% (+57.2 pp)} \\
Pooled False Alerts / 100 km$^2$ & 125.0 & 36.2 & \textbf{11.0 alerts / 100km$^2$} \\
Macro False Alerts / 100 km$^2$ & 128.3 & 30.2 & \textbf{8.2 alerts / 100km$^2$} \\
Negative Control Zone Violations & 5 violations & 0 violations & \textbf{0 violations (100\% clean)} \\
\bottomrule
\end{tabularx}
\end{table}

Furthermore, across all four localization purity regimes in our sensitivity sweep (\textit{Relaxed} $\text{PoC}\ge 0.02$, \textit{Standard} $\text{PoC}\ge 0.15$, \textit{Moderate} $\text{PoC}\ge 0.20$, and \textit{Strict} $\text{PoC}\ge 0.30$), **TerraShift v6.0 Hybrid** and **TerraShift v5.1** maintain Ranks \#1 and \#2 respectively, outperforming standalone AlphaEarth V1 by more than $+40$ percentage points of F1 in every regime.

\section{Hardware Efficiency and Resource Profiling}
\label{sec:res_hardware_profile}
Table \ref{tab:hardware_benchmarks} summarizes CPU latency and memory utilization across all seven evaluated models.

\begin{table}[htbp]
\centering
\caption{Hardware Latency and Memory Profiling on Commodity CPU.}
\label{tab:hardware_benchmarks}
\begin{tabularx}{\textwidth}{@{}lXXXX@{}}
\toprule
\textbf{Model Architecture} & \textbf{Inference (s)} & \textbf{Vectorization (s)} & \textbf{Pure Compute (s)} & \textbf{Peak RAM (MB)} \\
\midrule
\textbf{TerraShift v6.0 Hybrid} & \textbf{1.192 s} & \textbf{0.293 s} & \textbf{1.485 s} & \textbf{301 MB} \\
\textbf{TerraShift v5.1 (Siamese+L2A)} & \textbf{1.168 s} & \textbf{0.282 s} & \textbf{1.450 s} & \textbf{298 MB} \\
Google AlphaEarth V1 Cosine (A) & 0.210 s & 0.288 s & 0.498 s & 240 MB \\
Pure Multispectral Physical & 0.095 s & 0.461 s & 0.556 s & 185 MB \\
Untuned Siamese U-Net (RGB) & 0.354 s & 0.164 s & 0.518 s & 295 MB \\
Google Dynamic World V1 LULC & 0.215 s & 0.282 s & 0.497 s & 250 MB \\
AlphaEarth + Dynamic World (B) & 0.212 s & 0.252 s & 0.464 s & 265 MB \\
\bottomrule
\end{tabularx}
\end{table}

Pure CPU compute time for both TerraShift v6.0 Hybrid (**1.485 seconds**) and TerraShift v5.1 (**1.450 seconds**) is strictly under 1.5 seconds per corridor, with a peak memory footprint of **301 MB RAM**, satisfying both NFR-01 ($<2.0\text{ s}$) and NFR-02 ($<500\text{ MB}$).

\section{Error Analysis and Methodological Provenance Disclosure}
\label{sec:res_limitations}
In accordance with strict scientific transparency, we document the primary sources of residual error and reference label provenance:
\begin{enumerate}
    \item \textbf{Reference Parcel Provenance:} The 19 positive change events and 4 negative control zones in the Saudi pilot benchmark are coarse parcel-level bounding boxes curated from multi-date Sentinel-2 visual/spectral inspection and known Saudi urban and agricultural landmarks rather than surveyed municipal cadastral polygons. While our Localization-Aware 1-to-1 matching protocol ($\text{PoC} \ge 0.15, \text{IoU} \ge 0.05$) rigorously rejects sprawling detections, future operational qualification with MOMRAH and ROSHN should evaluate against dense surveyed GIS cadastres.
    \item \textbf{Duplicate Parcel Fragmentation ($FP_{\text{dup}}$):} Out of the 10 false positives generated by TerraShift v6.0 Hybrid across $91.19\text{ km}^2$, 4 are duplicate sub-parcel detections ($FP_{\text{dup}}$) where a large construction corridor or pivot cluster is split by an ungraded internal road, and only 6 are unassociated background alerts ($FP_{\text{bg}}$, yielding **76.0\% Pooled Validity Precision**).
    \item \textbf{Operational Cadence Complementarity:} Because Google AlphaEarth V1 embeddings are published annually, TerraShift v6.0 Hybrid is designed for annual strategic compliance audits, whereas TerraShift v5.1 operates on 5-day Sentinel-2 L2A acquisitions to provide continuous intra-year tactical monitoring.
\end{enumerate}
\newpage
